\section{中小企業退職金共済}

\subsection{加入対象企業・加入対象従業員}

\begin{fillblank}[年金1 2019 問題1(4)\textcircled{3}]
中小企業退職金共済に加入できる卸売業の企業の条件は、「常用従業員数（　　A　　）人以
下または資本金・出資金（　　B　　）円以下」となっている。
\end{fillblank}

\begin{answerbox}
  \begin{enumerate}[A:]
    \item 100
    \item 1億
  \end{enumerate}
\end{answerbox}

\begin{fillblank}[年金1 H28 問題1(3)]
  中小企業退職金共済について、次のA〜Eの空欄に当てはまる適切な語句を記入しな
さい。

  \begin{itemize}
    \item 加入できる従業員（被共済者）について、原則として事業主は雇用する従業員を全員加入さ
せなければならないが、
    \begin{itemize}
      \item 期間を定めて雇用される者
      \item 季節的業務に雇用される者
      \item （　　A　　）
      \item 短時間労働者
      \item （　　B　　）
      \item 相当の期間内に雇用関係の終了することが明らかな者
      \item 被共済者となることに反対する意思を表明した者
    \end{itemize}
といった条件にあてはまる者については、加入させなくてもよい。

    \item 新規加入助成と月額変更助成という 2 種類の掛金助成制度がある。
    \item 新規加入する事業主に対しては、国から掛金月額の （　　C　　） （掛金月額が 4,000 円以下の場合
    は加算あり。また、従業員ごとに上限 （　　D　　） 。 ）の助成がある。
    \item 18,000 円以下の掛金月額を増額変更する事業主に対しては、国から（　　E　　）の助成がある。
  \end{itemize}

\end{fillblank}

\begin{answerbox}
  \begin{enumerate}[A:]
    \item 試みの雇用期間中の者
    \item 休職期間中の者
    \item 2分の1
    \item 5,000円
    \item 増額分の3分の1
  \end{enumerate}
\end{answerbox}


\begin{fillblank}[数理人試験 H29 問題1 設問6]
  \noindent \textbf{◯中退共法}
  \begin{lawstructure}
    \begin{lawarticle}{第二条}{定義}
      この法律で「中小企業者」とは、次の各号のいずれかに該当する事業主（国、地方公共団体その他厚生労働省令で定めるこれらに準ずる者を除く。）をいう。

      \lawitem{一}
      常時雇用する従業員の数が（　　A　　）の事業主及び資本金の額又は出資の総額が（　　B　　）の法人である事業主（次号から第四号までに掲げる業種に属する事業を主たる事業として営む事業主を除く。）
      \lawitem{二}
      卸売業に属する事業を主たる事業として営む事業主であつて、常時雇用する従業員の数が（　　C　　）のもの及び資本金の額又は出資の総額が一億円以下の法人であるもの
      \lawitem{三}
      サービス業に属する事業を主たる事業として営む事業主であつて、常時雇用する従業員の数が（　　C　　）のもの及び資本金の額又は出資の総額が五千万円以下の法人であるもの
      \lawitem{四}
      （　　D　　）に属する事業を主たる事業として営む事業主であつて、常時雇用する従業員の数が五十人以下のもの及び資本金の額又は出資の総額が五千万円以下の法人であるもの

      \lawparagraph{2}
      この法律で「退職」とは、従業員について、事業主との雇用関係が終了することをいう。
      \lawparagraph{3}
      この法律で「退職金共済契約」とは、事業主が独立行政法人勤労者退職金共済機構（第五十六条及び第五十七条を除き、以下「機構」という。）に掛金を納付することを約し、機構がその事業主の雇用する従業員の退職について、この法律の定めるところにより、退職金を支給することを約する契約であつて、特定業種退職金共済契約以外のものをいう。
      \lawparagraph{4}
      この法律で「特定業種」とは、（　　E　　）その他従業員の相当数が、通常、当該業種に属する多数の事業の間を移動してこれらの事業の事業主に雇用される業種であつて、厚生労働大臣が指定するものをいう。
      \lawparagraph{5}
    この法律で「特定業種退職金共済契約」とは、特定業種に属する事業の事業主が機構に掛金を納付することを約し、機構が、期間を定めて雇用される者としてその事業主に雇用され、かつ、当該特定業種に属する事業に従事することを常態とする者の退職について、この法律の定めるところにより、退職金を支給することを約する契約をいう。
      \lawparagraph{6}
    この法律で「共済契約者」とは、退職金共済契約又は特定業種退職金共済契約の当事者である事業主をいう。
      \lawparagraph{7}
    この法律で「（　　F　　）」とは、退職金共済契約又は特定業種退職金共済契約により機構がその者の退職について退職金を支給すべき者をいう。
    \end{lawarticle}
  \end{lawstructure}
\end{fillblank}

\begin{answerbox}
  \begin{enumerate}[A:]
    \item 300人以下
    \item 3億円以下
    \item 100人以下
    \item 小売業
    \item 建設業
    \item 被共済者
  \end{enumerate}
\end{answerbox}

\newpage

\subsection{掛金・助成金}

\begin{fillblank}[数理人試験 2023 問題1 設問5]
  次は、中小企業退職金共済制度の掛金に関する記述である。

  \begin{itemize}
  \item 短時間労働者を除く被共済者の中小企業退職金共済の掛金は、月額が（　　A　　）円から（　　B　　）円までの1,000円刻みと12,000円から（　　C　　）円までの2,000円刻みの中から被共済者ごとに選択でき、事業主が全額負担する。

  \item 新規加入助成と月額変更助成という2種類の掛金助成制度がある。新しく加入する事業主に対しては、加入後4か月目から1年間、国から掛金月額の（　　D　　）の助成がある。

    また、（　　E　　）円以下の掛金月額を増額変更する事業主に対して、増額する月から1年間、国から増額分（増額前と増額後の掛金月額の差額）の（　　F　　）の助成がある。
\end{itemize}
  
\end{fillblank}

\begin{answerbox}
  \begin{enumerate}[A:]
    \item 5,000
    \item 10,000
    \item 30,000
    \item 2分の1
    \item 18,000
    \item 3分の1
  \end{enumerate}
\end{answerbox}


\begin{fillblank}[年金1 2018 問題1(5)\textcircled{5}]
中小企業退職金共済の被共済者（短時間労働者を除く）1 人あたりの掛金月額の下限は
（　　A　　）円となっている。
\end{fillblank}

\begin{answerbox}
  \begin{enumerate}[A:]
    \item 5,000
  \end{enumerate}
\end{answerbox}

\begin{shortanswer}[年金1 2023 問題2(4)ア]
新しく中小企業退職金共済制度に加入する事業主に対する、国の助成について簡潔に入力し
なさい。 （250 字以内）
\end{shortanswer}

\begin{answerbox}
  以下の金額を国が助成する。
  \begin{enumerate}[(1)]
    \item 加入後4ヶ月目から1年間について、掛金月額の2分の1（従業員ごとに上限5,000円）
    \item 短時間労働者の特例掛金月額（掛金月額4,000円以下）加入者については(1)に次の額を上乗せする
    
    掛金月額2000円の場合：300円

    掛金月額3000円の場合：400円

    掛金月額4000円の場合：500円
  \end{enumerate}
  
\end{answerbox}

\begin{shortanswer}[年金1 2023 問題2(4)イ]
掛金月額を増額する事業主に対する、国の助成について簡潔に入力しなさい。 （250 字以内）
\end{shortanswer}

\begin{answerbox}
  18,000円以下の掛金月額を拠出している従業員の掛金を増額する事業主に対し、増額月から1年間について、増額分の3分の1を国が助成する。
\end{answerbox}


\begin{fillblank}[年金1 2024 問題1(6)イ]
中小企業退職金共済の掛金月額の下限は5,000 円（パートタイマー等短時間労働者は2,000 円）、
上限は（　　A　　）円である。
\end{fillblank}

\begin{answerbox}
  \begin{enumerate}[A:]
    \item 30,000
  \end{enumerate}
\end{answerbox}


\begin{fillblank}[数理人試験 2022 問題1 設問5]
  \noindent \textbf{◯中退共法}
  \begin{lawstructure}
    \begin{lawarticle}{第二十二条}{掛金の納付}
      共済契約者は、退職金共済契約が効力を生じた日の（　　A　　）から被共済者が退職した日又は退職金共済契約が解除された日の（　　B　　）までの各月につき、その月の末日（退職の日又は退職金共済契約の解除の日の（　　B　　）にあつては、その退職の日又はその解除の日）における掛金月額により、毎月分の掛金を翌月末日（退職金共済契約が効力を生じた日の（　　A　　）分の掛金にあつては、翌々月末日）までに納付しなければならない。

      \lawparagraph{2}
      毎月分の掛金は、（　　C　　）することができない。
    \end{lawarticle}


    \begin{lawarticle}{第二十三条}{加入促進等のための掛金負担軽減措置}
      機構は、中小企業者が退職金共済契約の申込みをすること及び共済契約者が第九条第一項の（　　D　　）の申込みをすることを促進するため、厚生労働省令で定めるところにより、共済契約者の掛金に係る負担を軽減する措置として、（　　E　　）の掛金の額を減額することができる。

      \lawparagraph{2}
      前項の規定に基づき掛金の減額の措置が講ぜられる月について、共済契約者が同項の規定に基づき減額された額により掛金を納付した場合には、第十条第二項（第十六条第三項において準用する場合を含む。）及び第五十五条第一項の規定の適用については、前条第一項の掛金月額により掛金の納付があつたものとみなす。
    \end{lawarticle}

      \begin{lawarticle}{第二十四条}{}
        （略）
    \end{lawarticle}
    
    \begin{lawarticle}{第二十五条}{（　　F　　）}
機構は、納付期限後に掛金を納付する共済契約者に対して、（　　F　　）を納付させることができる。

      \lawparagraph{2}
（　　F　　）の額は、掛金の額につき年十四・六パーセントの割合で納付期限の翌日から納付の日の前日までの日数によつて計算した額を超えてはならない。
    \end{lawarticle}

  \end{lawstructure}
\end{fillblank}

\begin{answerbox}
  \begin{enumerate}[A:]
    \item 属する月
    \item 属する月
    \item 分割して納付
    \item 掛金月額の増加
    \item 一定の月分
    \item 割増金
  \end{enumerate}
\end{answerbox}

\newpage

\subsection{退職金}


\begin{fillblank}[年金1 2021 問題1(6)\textcircled{1}（改題）]
中小企業退職金共済の付加退職金は、運用利回りが予定利回りを上回った場合に基本退職金
に上積みされるものである。具体的には、掛金納付月数の（　　A　　）か月目とその後12 か月ご
との基本退職金相当額に、厚生労働大臣が定めるその年度の支給率を乗じて得た額を、退職時
まで累計した金額となり、\warekiyear 年度の支給率は（　　B　　）となる。
\end{fillblank}

\begin{answerbox}
  \begin{enumerate}[A:]
    \item 43
    \item 0.0061 %毎年修正が必要
  \end{enumerate}
\end{answerbox}

\begin{explanationbox}
Bは毎年更新改められる。

\url{https://chutaikyo.taisyokukin.go.jp/notice/new/osirase02.html}
\end{explanationbox}

\begin{fillblank}[年金1 2020 問題1(6)]
中小企業退職金共済の給付について、基本退職金と付加退職金を合算した退職金の額の全額
を5 年間または10 年間にわたって分割して受け取ることができる要件は次のとおりである。
\begin{itemize}
  \item 退職日に（　　A　　）歳以上であること
\item 退職金の額が、5 年間の分割払いの場合は（　　B　　）万円以上であり、10 年間の分割払いの
場合は150 万円以上であること
\end{itemize}

\end{fillblank}


\begin{answerbox}
  \begin{enumerate}[A:]
    \item 60
    \item 80
  \end{enumerate}
\end{answerbox}


\begin{fillblank}[数理人試験 2020 問題1 設問6 A〜C]
    \noindent \textbf{◯中退共法}
  \begin{lawstructure}
    \begin{lawarticle}{第十条}{退職金}
      機構は、被共済者が退職したときは、その者（退職が死亡によるものであるときは、その遺族）に退職金を支給する。ただし、当該被共済者に係る掛金の納付があつた月数（以下「掛金納付月数」という。）が（　　A　　）に満たないときは、この限りでない。

      \lawparagraph{2}
退職金の額は、次の各号に掲げる掛金納付月数の区分に応じ、当該各号に定める額とする。

      \lawitem{一}
二十三月以下　被共済者に係る（　　B　　）を下回る額として、（　　C　　）及び掛金納付月数に応じ政令で定める額（退職が死亡による場合にあつては、被共済者に係る（　　B　　）に相当する額として、（　　C　　）及び掛金納付月数に応じ政令で定める額）

      \lawitem{二}
二十四月以上四十二月以下　被共済者に係る（　　B　　）に相当する額として、（　　C　　）及び掛金納付月数に応じ政令で定める額

      \lawitem{三}
四十三月以上　次のイ及びロに定める額を合算して得た額

      \hangedlawitem{イ}{3}
被共済者に係る（　　B　　）を上回る額として、（　　C　　）及び掛金納付月数に応じ政令で定める額

      \hangedlawitem{ロ}{3}
計算月（その月分の掛金の納付があつた場合に掛金納付月数が四十三月又は四十三月に十二月の整数倍の月数を加えた月数となる月をいう。以下この号及び第四項において同じ。）に被共済者が退職したものとみなしてイの規定を適用した場合に得られる額（第四項において「仮定退職金額」という。）に、それぞれ当該計算月の属する年度（四月一日から翌年三月三十一日までをいう。以下同じ。）に係る支給率を乗じて得た額（その額に一円未満の端数があるときは、これを一円に切り上げるものとする。）を合算して得た額

（以下略）
    \end{lawarticle}
  \end{lawstructure}
\end{fillblank}

\begin{answerbox}
  \begin{enumerate}[A:]
    \item 十二月
    \item 納付された掛金総額
    \item 掛金月額
  \end{enumerate}
\end{answerbox}

\begin{fillblank}[数理人試験 2020 問題1 設問6 D,E]
  \noindent \textbf{◯中退共法}
  \begin{lawstructure}
    \begin{lawarticle}{第十二条}{退職金の分割支給等}
      機構は、前条の規定にかかわらず、被共済者の請求により、退職金の全部又は一部を分割払の方法により支給することができる。ただし、次の各号のいずれかに該当する場合は、この限りでない。

      \lawitem{一}
      退職金の額が厚生労働省令で定める金額未満であるとき。

      \lawitem{二}
被共済者が退職した日において（　　A　　）であるとき。

\lawitem{三}
被共済者が退職金の一部を分割払の方法により支給することを請求した場合において、次項に規定する分割払対象額が厚生労働省令で定める金額未満であるとき又は当該退職金の全額から同項に規定する分割払対象額を減じた額が厚生労働省令で定める金額未満であるとき。

\lawparagraph{2}
被共済者が退職金の一部について分割払の方法により支給を受けようとする場合における前項の請求は、当該分割払の方法により支給を受けようとする退職金の一部の額（以下この条において「分割払対象額」という。）を定めてしなければならない。

\lawparagraph{3}
分割払の方法による退職金の支給期月は、（　　B　　）とする。

（以下略）
    \end{lawarticle}
  \end{lawstructure}
\end{fillblank}

\begin{answerbox}
  \begin{enumerate}[A:]
    \item 六十歳未満
    \item 毎年二月、五月、八月及び十二月
  \end{enumerate}
\end{answerbox}


\begin{fillblank}[数理人試験 H28 問題1 設問6]
  \noindent \textbf{◯中退共法}
  \begin{lawstructure}
    \begin{lawarticle}{第十七条}{}
      第八条第二項第二号の規定により退職金共済契約が（　　A　　）された際に、当該（　　A　　）された退職金共済契約の共済契約者が、
      当該（　　A　　）された退職金共済契約の被共済者に係る確定給付企業年金法（平成十三年法律第五十号）第二条第一項に規定する
      確定給付企業年金（第三十一条の三及び第三十一条の四において「確定給付企業年金」という。）、
      確定拠出年金法（平成十三年法律第八十八号）第二条第二項に規定する（　　B　　）
      （第三十一条の三及び第三十一条の四において「（　　B　　）」という。）その他の政令で定める制度であつて、
      厚生労働省令で定める要件を備えているもの（以下この条において「特定企業年金制度等」という。）
      の（　　C　　）の通知をした場合には、前条第一項の規定にかかわらず、機構は、当該被共済者に（　　D　　）を支給しない。
      この場合において、当該共済契約者が、当該（　　A　　）後厚生労働省令で定める期間内に、当該被共済者の同意を得て、
      厚生労働省令で定めるところにより、当該通知に係る特定企業年金制度等への（　　D　　）に相当する額の引渡しに関する
      申出をしたときは、機構は、当該申出に基づき、当該被共済者に係る（　　D　　）に相当する額の範囲内の金額で厚生労働省令
      で定める金額を、確定給付企業年金法第三十条第三項に規定する資産管理運用機関等（第三十一条の三及び第三十一条の四
      において「資産管理運用機関等」という。）、確定拠出年金法第二条第七項第一号ロに規定する資産管理機関
      （第三十一条の三及び第三十一条の四において「資産管理機関」という。）その他の当該特定企業年金制度等を実施する団体
      として厚生労働省令で定めるものに引き渡すものとする。

      （以下略）
    \end{lawarticle}

    \begin{lawarticle}{第三十一条の四}{産管理運用機関等への解約手当金に相当する額の移換等}
      共済契約者が（　　E　　）（平成十七年法律第八十六号）その他の法律の規定による合併、
      （　　F　　）その他の行為として厚生労働省令で定める行為（以下この項において「合併等」という。）をした場合であつて、
      当該合併等により退職金共済契約が第八条第三項第一号の規定に基づき（　　A　　）された被共済者を加入者とする確定給付企業年金
      又は（　　B　　）加入者とする（　　B　　）を実施するときは、機構は、当該共済契約者が当該被共済者の同意を得て厚生労働省令
      で定めるところにより行う確定給付企業年金又は（　　B　　）（厚生労働省令で定めるものに限る。）への
      （　　D　　）に相当する額の移換に関する申出に基づき、資産管理運用機関等又は資産管理機関に当該同意を得た
      被共済者に係る（　　D　　）に相当する額を移換するものとする。

      （以下略）
    \end{lawarticle}
  \end{lawstructure}
\end{fillblank}

\begin{answerbox}
  \begin{enumerate}[A:]
    \item 解除
    \item 企業型年金
    \item 実施
    \item 解約手当金
    \item 会社法
    \item 会社分割
  \end{enumerate}
\end{answerbox}

\newpage

\subsection{他制度との通算・資産移換}

\begin{shortanswer}[年金1 2020 問題2(5)]
中小企業退職金共済制度の通算制度に関し、次の(1)および(2)の設問に答えなさい。
\begin{enumerate}[(1)]
  \item 過去勤務期間の通算が可能となる対象者、要件および通算限度期間を簡記しなさい。
\item 転職により掛金納付実績の通算が可能となる他の共済制度を挙げなさい。
\end{enumerate}
\end{shortanswer}

\begin{answerbox}
  \begin{enumerate}[(1)]
    \item 新規加入する際、1年以上勤務している従業員全員について、通常の掛金とは別に過去勤務掛金を納付することにより、10年を限度として加入前の勤務期間の通算が可能。
    \item 特定業種退職金共済制度または特定退職金共済制度
  \end{enumerate}
\end{answerbox}

\begin{fillblank}[数理人試験 H30 問題1 設問6]
    \noindent \textbf{◯中退共則}
  \begin{lawstructure}
    \begin{lawarticle}{第三十一条}{法第十七条第一項の厚生労働省令で定める要件}
      法第十七条第一項の厚生労働省令で定める要件は、次の各号に掲げる制度の区分に応じ、当該各号に定めるものとする。

      \lawitem{一}
中小企業退職金共済法施行令（昭和三十九年政令第百八十八号。以下「令」という。）第三条第一号の確定給付企業年金（以下「確定給付企業年金」という。）　次のイからハまでのいずれにも該当すること。

      \hangedlawitem{イ}{3}
法第八条第二項第二号の規定により（　　A　　）された共済契約の被共済者の（　　B　　）を確定給付企業年金法（平成十三年法律第五十号）第二条第四項に規定する（　　C　　）（以下「（　　C　　）」という。）とするものであること。

      \hangedlawitem{ロ}{3}
法第十七条第一項の引渡しをしたときにおける同項後段の申出に係る被共済者に係る確定給付企業年金法施行規則（平成十四年厚生労働省令第二十二号）第四十三条の規定に基づき計算した
（　　D　　）から当該引渡しがないものとして同条の規定に基づき計算した
（　　D　　）を控除した額が、当該被共済者に係る第三十五条に規定する金額の合算額を（　　E　　）ものであること。

\hangedlawitem{ハ}{3}
法第十七条第一項の規定により機構が引き渡す金額が、
同項後段の申出をする共済契約者が負担する掛金として（　　E　　）払い込まれるものであること。
    
（以下略）
\end{lawarticle}
  \end{lawstructure}
\end{fillblank}

\begin{answerbox}
  \begin{enumerate}[A:]
    \item 解除
    \item 全て
    \item 加入者
    \item 給付に要する費用の予測に基づく予測がくの現価に相当する額
    \item 下回らない
    \item 一括して
  \end{enumerate}
\end{answerbox}

\begin{shortanswer}[年金1 2025 問題2(4)]
  中小企業退職金共済制度に関し、次の設問に解答しなさい。 
\\

中小企業でなくなったことにより退職金共済契約が解除された際に生じた解約手当金につい
て、一定の要件を満たせば確定給付企業年金制度に移換することができる。この解約手当金相
当額を確定給付企業年金制度に移換するための一定の要件として、中小企業退職金共済法施行
規則第３１条の第１号に規定されている内容を簡潔に入力しなさい。 （250 字以内）
\\

（参考）

中小企業退職金共済法施行規則第３１条には、解約手当金に相当する額の引き渡し先となる制
度に応じて第１号から第３号（第１号：確定給付企業年金、第２号：企業型年金、第３号：特
定退職金共済制度）までの規定があり、各号には各制度に解約手当金に相当する額を引き渡す
ための一定の要件がそれぞれ定められている。
\end{shortanswer}

\begin{answerbox}
  次のAからCまでのいずれにも該当すること。

  \begin{enumerate}[A:]
    \item 解除された共済契約の被共済者の全てを確定給付企業年金の加入者とするものであること。
    \item 解約手当金の引渡しの対象となる被共済者に係る通常予測給付現価から当該引渡しがない
ものとして計算した通常予測給付現価を控除した額が、当該被共済者に係る解約手当金の合
算額を下回らないものであること。
\item 解約手当金として機構が引き渡す金額が、引き渡しの申出をする共済契約者が負担する掛
金として一括して払い込まれるものであること。
  \end{enumerate}

\end{answerbox}

\newpage